\begin{problem}[2020第37届全国中学生物理竞赛复赛][高铁空气弹簧减振系统]{117}
    高铁运行的平稳性与列车的振动程度密切相关，列车上安装的空气弹簧可以有效减振。某高铁测试实验采用的空气弹簧模型由主气室（气囊）和附加气室（容积不变）构成。空气弹簧对簧上负载竖直向上的作用力由气囊内被压缩的空气产生的弹力提供。簧上负载处于平衡状态时，主气室内气体的压强和体积分别为 $p_{10}$ 和 $V_{10}$，附加气室的容积为 $V_{2}$，大气压强为 $p_{0}$。空气弹簧竖直向上的作用力 $F$ 与主气室内气体压强和大气压强之差的比值称为空气弹簧的有效承载面积 $A_{\mathrm{e}}$。已知空气的定容摩尔热容 $C_{V} = \frac{5}{2} R$（$R$ 是普适气体常量）。假设在微振幅条件下主气室内气体的体积 $V_{1}$ 和有效承载面积 $A_{\mathrm{e}}$ 均可视为空气弹簧（气囊）的承载面相对于其平衡位置的竖直位移 $y$（向下为正）的线性函数，变化率分别为常量 $-\alpha (\alpha > 0)$ 和 $\beta (\beta > 0)$。
    \begin{subproblem}
        附加气室未与主气室连通（阀门关闭），试在下述两种情形下，导出空气弹簧的劲度系数 $K \equiv \mathrm{d} F / \mathrm{d} y$ 与其有效承载面积 $A_{\mathrm{e}}$ 之间的关系：
        \begin{subsubproblem}
            上下乘客（主气室内气体压强和体积的变化满足等温过程）；
        \end{subsubproblem}
        \begin{subsubproblem}
            列车运行中遇到剧烈颠簸（主气室内气体压强和体积变化满足绝热过程）。
        \end{subsubproblem}
    \end{subproblem}
    \begin{subproblem}
        主气室连通附加气室（阀门打开）后，在上述两种情形下，导出空气弹簧的劲度系数 $K$ 与其有效承载面积 $A_{\mathrm{e}}$ 之间的关系。主气室与附加气室之间连通管道的容积可忽略。
    \end{subproblem}
\end{problem}